\section{开放式评估（Open-ended Evaluation）}

\subsection{定义与挑战}

开放式任务与封闭式任务相反：存在\textbf{大量可能的正确答案}，且无法穷举所有答案。更重要的是，答案质量是一个\textbf{连续谱}——不是简单的对/错，而是有好有坏的程度之分。

典型的开放式任务包括：
\begin{itemize}
    \item \textbf{文本摘要}（Summarization）：将长文本压缩为简短摘要。典型基准：CNN/Daily Mail
    \item \textbf{机器翻译}（Translation）：在不同语言间翻译文本
    \item \textbf{指令跟随}（Instruction Following）：类似 ChatGPT 的对话，这是所有任务的``母任务''——任何之前的任务都可以被视为对聊天机器人的一条指令
\end{itemize}

\begin{figure}[H]
\centering
\includegraphics[width=0.95\textwidth]{slides-images/slide-015.jpg}
\caption{文本生成评估的三大类方法：内容重叠指标（Content Overlap Metrics）、基于模型的指标（Model-based Metrics）和人工评估（Human Evaluations）\protect\footnotemark}
\end{figure}
\footnotetext{Slides 第16页。部分 slides 改编自 Asli Celikyilmaz 在 EMNLP 2020 tutorial。}

\subsection{内容重叠指标（Content Overlap Metrics）}

最简单的评估方式是将模型生成的文本与人工撰写的参考答案逐词比较。

\begin{importantbox}{BLEU 和 ROUGE}
\begin{itemize}
    \item \textbf{BLEU}（Bilingual Evaluation Understudy）：关注\textbf{精确率}——生成文本中有多少 n-gram 出现在参考文本中。主要用于机器翻译。还包含长度惩罚，防止只生成``the''等高频词来获得高精确率
    \item \textbf{ROUGE}（Recall-Oriented Understudy for Gisting Evaluation）：关注\textbf{召回率}——参考文本中有多少 n-gram 出现在生成文本中。主要用于文本摘要
\end{itemize}
这两个指标直到约 2 年前还是翻译和摘要任务的``金标准''。
\end{importantbox}

\begin{figure}[H]
\centering
\includegraphics[width=0.95\textwidth]{slides-images/slide-016.jpg}
\caption{N-gram 重叠指标示意：比较生成序列（红色标注）与参考序列中共有的词\protect\footnotemark}
\end{figure}
\footnotetext{Slides 第17页。}

\subsubsection{内容重叠指标的局限性}

讲者用一个生动的例子说明了这些指标的缺陷。假设问题是``你喜欢 CS224N 的课吗？''，参考答案是``Heck yes!''：

\begin{figure}[H]
\centering
\includegraphics[width=0.95\textwidth]{slides-images/slide-018.jpg}
\caption{BLEU 分数的缺陷示例：``Yes'' 得 67\%（假正例），``Yep'' 得 0\%（假负例，但语义完全相同），``Heck no'' 得 67\%（假正例，语义完全相反）\protect\footnotemark}
\end{figure}
\footnotetext{Slides 第19页。}

\begin{warningbox}{词汇重叠 $\neq$ 语义相似}
\begin{itemize}
    \item \textbf{假负例}：``Yep'' 与 ``Heck yes!'' 语义完全相同，但 BLEU = 0\%，因为没有词重叠
    \item \textbf{假正例}：``Heck no'' 与 ``Heck yes!'' 语义完全相反，但 BLEU = 67\%，因为大部分词相同
\end{itemize}
内容重叠指标根本无法捕捉语义——它们只看表面的词汇匹配。
\end{warningbox}

\subsection{基于嵌入的指标}

既然词级匹配无法捕捉语义，自然的改进方向是使用\textbf{词嵌入}（word embeddings）或\textbf{上下文嵌入}（contextual embeddings）来比较语义相似度。

\begin{figure}[H]
\centering
\includegraphics[width=0.95\textwidth]{slides-images/slide-020.jpg}
\caption{基于嵌入的评估：使用词嵌入的平均向量来计算语义相似度，以及使用 BERT 产生的上下文嵌入（BERTScore）\protect\footnotemark}
\end{figure}
\footnotetext{Slides 第21页。}

\begin{knowledgebox}{从词嵌入到 BERTScore}
\begin{itemize}
    \item \textbf{词嵌入平均}（约 2016 年）：将参考序列和生成序列的词嵌入分别取平均，再计算余弦相似度。优点：简单。缺点：丢失了上下文信息
    \item \textbf{BERTScore}（约 2019 年）：将两个序列分别通过 BERT 获得上下文嵌入，再进行``智能匹配''。BERTScore 至今仍被广泛使用，效果显著优于 BLEU/ROUGE
    \item \textbf{BLEURT}：在 BERT 基础上，先用 BLEU 等指标作为无监督预训练目标，再用人工标注数据微调，使模型学会像人类一样评估文本质量
\end{itemize}
\end{knowledgebox}

\subsection{参考答案质量的影响}

\begin{figure}[H]
\centering
\includegraphics[width=0.95\textwidth]{slides-images/slide-023.jpg}
\caption{参考答案质量对 ROUGE-L 有效性的影响：使用标准参考时（左），ROUGE-L 与人工评估几乎不相关；使用专家撰写的高质量参考时（右），相关性显著提高\protect\footnotemark}
\end{figure}
\footnotetext{Slides 第24页。}

\begin{importantbox}{参考答案是瓶颈，不只是指标}
上图展示了一个关键发现：ROUGE-L 与人工评估的相关性之所以低，很大程度上是因为参考答案质量差（如 CNN 文章顶部的简略要点），而非 ROUGE 本身完全无用。当使用专家撰写的高质量参考时，相关性显著提高。这说明\textbf{基于参考的评估指标只能和参考答案一样好}。
\end{importantbox}

\subsection{本章小结}

开放式评估远比封闭式评估复杂。传统的内容重叠指标（BLEU、ROUGE）无法捕捉语义相似性，存在严重的假正例和假负例问题。基于嵌入的方法（BERTScore、BLEURT）有所改进但仍不完美。所有基于参考的方法都受制于参考答案的质量。这些局限性促使领域转向无参考评估和人工评估。

%% ========== Part 4: 人工评估 ==========

